% Source: https://artofproblemsolving.com/community/c3h2245401_super_binomial_theorem
% Stable identity: aops-topic-2245401-q01
\textbf{(a)} $(a_1+\cdots+a_k)^n$ is a product of $n$ factors, each equal to $a_1+\cdots+a_k$. Multiplying out, we choose one term $a_i$ from each factor and multiply the choices; the expansion is the sum of all $k^n$ such products. A choice gives $a_1^{m_1}\cdots a_k^{m_k}$ exactly when $a_1$ is chosen from $m_1$ of the factors, $a_2$ from $m_2$ of them, and so on. Such a choice is the same as a word of length $n$ (the $i$th letter tells which term is chosen from the $i$th factor) containing $m_1$ letters $a_1$, \ldots, $m_k$ letters $a_k$. The number of these words is
\[\binom{n}{m_1}\binom{n-m_1}{m_2}\cdots\binom{m_k}{m_k}=\dfrac{n!}{m_1!\,(n-m_1)!}\cdot\dfrac{(n-m_1)!}{m_2!\,(n-m_1-m_2)!}\cdots=\dfrac{n!}{m_1!\,m_2!\cdots m_k!}.\]
Collecting equal products gives the stated coefficient. For $k=2$ this is the binomial theorem, since $\dfrac{n!}{m!\,(n-m)!}=\dbinom nm$. $\blacksquare$

\textbf{(b)} The product index is $f$, but the factor $m_k!$ does not depend on $f$; as written, $\prod_{f=1}^k m_k!=(m_k!)^k$. The correct expression is $\prod_{f=1}^k m_f!=m_1!\,m_2!\cdots m_k!$.

\textbf{(c)} Here $n=7$ and the exponents are $2,1,3,1$, so the coefficient is
\[\dfrac{7!}{2!\,1!\,3!\,1!}=\dfrac{5040}{12}=\boxed{420}.\]
